% Generated draft: requires mathematical and source review.
\documentclass[11pt]{article}
\usepackage[margin=25mm]{geometry}
\usepackage{fontspec}
\setmainfont{Latin Modern Roman}
\usepackage{amsmath,amssymb,mathtools,unicode-math}
\setmathfont{Latin Modern Math}
\usepackage{enumitem,needspace,xurl,hyperref}
\setlength{\parindent}{0pt}
\setlength{\parskip}{5pt}
\setlength{\emergencystretch}{3em}
\newcommand{\N}{\mathbb{N}}
\newcommand{\Z}{\mathbb{Z}}
\newcommand{\Q}{\mathbb{Q}}
\newcommand{\R}{\mathbb{R}}
\newcommand{\C}{\mathbb{C}}
\newcommand{\F}{\mathbb{F}}
\newcommand{\A}{\mathbb{A}}
\newcommand{\PP}{\mathbb{P}}
\newcommand{\E}{\mathbb{E}}
\newcommand{\eps}{\varepsilon}
\newcommand{\dd}{\,\mathrm{d}}
\begin{document}
\section*{UnsolvedMath ID: 30000103}
\textbf{Source problem code:} catalogue entry
\subsection{Definitions and mathematical statement}
Let $X$ be a smooth projective Fano variety, so $-K_X$ is ample. In the usual complex, characteristic-zero interpretation, define the codimension-two class
\[
 \eta_X=c_2(\Omega_X^1)-\tfrac12c_1(K_X)^2.
\]
The verified question asks whether negativity of this class forces $X$ to be rationally $1$-connected. The source gives as a concrete negativity criterion
\[
 \eta_X\cdot[S]<0\quad\text{for every integral surface }S\subset X.
\]
Here the product is a numerical degree on the surface. Since $c_2(\Omega_X^1)=c_2(T_X)$, this criterion is equivalently strict positivity of
\[
 \operatorname{ch}_2(T_X)=\tfrac12c_1(T_X)^2-c_2(T_X)=-\eta_X
\]
on every such surface. No uniform lower bound or separate ampleness condition on a codimension-two class is added.

A conventional rational-$1$-connectedness formulation uses an effective curve class $\beta$ and an irreducible family $M\subset\overline{\mathcal M}_{0,2}(X,\beta)$ of genus-zero stable maps with two marked points. Its evaluation map
\[
 \operatorname{ev}:M\longrightarrow X\times X,\qquad
 [f:C\to X,p_1,p_2]\longmapsto(f(p_1),f(p_2))
\]
should be dominant and have rationally connected geometric generic fiber. Rational connectedness means that two general points of that fiber can be joined by a rational curve. The question asks whether the Fano and stated Chern-negativity conditions guarantee the appropriate rational-$1$-connectedness property.

The extract does not specify the base field, the full meaning of ``negative'' beyond its illustrative surface-intersection test, or whether rational $1$-connectedness means the existence of one such component or a stronger assertion in all sufficiently positive degrees. The displayed component formulation is an explicit standard interpretation; those quantifiers must match the report before a stronger version is asserted.

The literature assessment checked on 21 August 2026 reports important low-degree and general-hypersurface results but no theorem resolving the general Chern-class criterion. Those special cases are not a proof for every smooth Fano variety satisfying the displayed inequality. The open label is dated triage requiring independent verification.
\subsection{Short English statement}
Does a smooth Fano variety with negative $c_2(\Omega_X^1)-\frac12K_X^2$, in particular negative intersection with every surface, have rationally connected families of rational curves through two general points? The exact rational-$1$-connectedness degree quantifiers need source confirmation.
\subsection{Sources}
\begin{itemize}
\item \emph{Arithmetic Algebraic Geometry} (2005), Oberwolfach Reports, DOI 10.4171/owr/2004/37, pp. 1971--2014, \url{https://doi.org/10.4171/owr/2004/37}; proposed year supplied as 2004.
\item Literature assessment checked 21 August 2026: original report, \url{https://doi.org/10.4171/owr/2004/37}; Starr, ``rational curves on low-degree hypersurfaces,'' \url{https://arxiv.org/abs/math/0305432}.
\item Local UnsolvedMath registry, permanent ID $30000103$, exact source code \texttt{catalogue entry}, created 25 August 2026 and updated 16 September 2026. Its clean statement is labelled corrected and verified, with the precise Chern-class hypothesis and illustrative surface-intersection criterion recovered from the report.
\end{itemize}
\end{document}
